\chapter{2010年重庆卷（理）}

\section{单选题}

\begin{problem}
在等比数列 \(\{a_{n}\}\) 中，\(a_{2010} = 8a_{2007}\)，则公比 \(q\) 的值为（\quad）

\choices
    {2}
    {3}
    {4}
    {8}
\end{problem}

\begin{problem}
已知向量 \(a, b\) 满足 \(a \cdot b = 0, |a| = 1, |b| = 2\)，则 \(|2a - b| =\)（\quad）

\choices
    {0}
    {\( 2\sqrt{2} \)}
    {4}
    {8}
\end{problem}

\begin{problem}
\(\lim_{x\to 2}\left(\frac{4}{x^2 - 4} -\frac{1}{x - 2}\right) =\)（\quad）

\choices
    {\(-1\)}
    {\( -\frac{1}{4} \)}
    {\( \frac{1}{4} \)}
    {1}
\end{problem}

\begin{problem}
设变量 \(x, y\) 满足约束条件 \(\left\{ \begin{array}{l} y \geqslant 0, \\ x - y + 1 \geqslant 0, \\ x + y - 3 \leqslant 0, \end{array} \right.\) 则 \(z = 2x + y\) 的最大值为（\quad）

\choices
    {\(-2\)}
    {4}
    {6}
    {8}
\end{problem}

\begin{problem}
函数 \( f(x) = \frac{4^x + 1}{2^x} \) 的图象（\quad）

\choices
    {关于原点对称}
    {关于直线 \( y = x \) 对称}
    {关于 \(x\) 轴对称}
    {关于 \(y\) 轴对称}
\end{problem}

\begin{problem}
已知函数 \( y = \sin (\omega x + \varphi) (\omega > 0, |\varphi| < \frac{\pi}{2}) \) 的部分图象如图所示. 则（\quad）

\bitmapfigure[max width=0.42\linewidth,max totalheight=4.8cm]{img/2010/chongqing_science/q6_fig1.png}

\choices
    {\(\omega = 1, \varphi = \frac{\pi}{6}\)}
    {\(\omega = 1, \varphi = -\frac{\pi}{6}\)}
    {\(\omega = 2, \varphi = \frac{\pi}{6}\)}
    {\(\omega = 2, \varphi = -\frac{\pi}{6}\)}
\end{problem}

\begin{problem}
已知 \(x > 0, y > 0, x + 2y + 2xy = 8\)，则 \(x + 2y\) 的最小值是（\quad）

\choices
    {3}
    {4}
    {\( \frac{9}{2} \)}
    {\(\frac{11}{2}\)}
\end{problem}

\begin{problem}
直线 \( y = \frac{\sqrt{3}}{3} x + \sqrt{2} \) 与圆心为 \( D \) 的圆 \( \left\{ \begin{array}{l} x = \sqrt{3} + \sqrt{3}\cos \theta, \\ y = 1 + \sqrt{3}\sin \theta, \end{array} \right. (\theta \in [0,2\pi)) \) 交于 \( A, B \) 两点，则直线 \( AD \) 与 \( BD \) 的倾斜角之和为（\quad）

\choices
    {\(\frac{7}{6}\pi\)}
    {\(\frac{5}{4}\pi\)}
    {\(\frac{4}{3}\pi\)}
    {\(\frac{5}{3}\pi\)}
\end{problem}

\begin{problem}
某单位安排 7 位员工在 10 月 1 日至 7 日值班，每天 1 人，每人值班 1 天，若 7 位员工中的甲，乙排在相邻两天. 丙不排在 10 月 1 日，丁不排在 10 月 7 日，则不同的安排方案共有（\quad）

\choices
    {504 种}
    {960 种}
    {1008 种}
    {1108 种}
\end{problem}

\begin{problem}
到两互相垂直的异面直线的距离相等的点，在过其中一条直线且平行于另一条直线的平面内的轨迹是（\quad）

\choices
    {直线}
    {椭圆}
    {抛物线}
    {双曲线}
\end{problem}

\section{填空题}

\begin{problem}
已知复数 \(z = 1 + \i\)，则 \(\frac{2}{z} - z =\fillinblank{}\).
\end{problem}

\begin{problem}
设 \(U = \{0, 1, 2, 3\}\)，\(A = \{x \in U \mid x^2 + mx = 0\}\)，若 \(\complement_U A = \{1, 2\}\)，则实数 \(m =\fillinblank{}\).
\end{problem}

\begin{problem}
某篮球队员在比赛中每次罚球的命中率相同，且在两次罚球中至多命中一次的概率为 \(\frac{16}{25}\)，则该队员每次罚球的命中率为\fillinblank{}.
\end{problem}

\begin{problem}
已知以 \(F\) 为焦点的抛物线 \(y^{2} = 4x\) 上的两点 \(A\)，\(B\) 满足 \(\overrightarrow{AF} = 3\overrightarrow{FB}\)，则弦 \(AB\) 的中点到准线的距离为\fillinblank{}.
\end{problem}

\begin{problem}
已知函数 \( f(x) \) 满足: \( f(1) = \frac{1}{4}, 4f(x)f(y) = f(x + y) + f(x - y) (x, y \in \R) \)，则 \( f(2010) = \)\fillinblank{}.
\end{problem}

\section{解答题}

\begin{problem}
设函数 \( f(x) = \cos \left(x + \frac{2\pi}{3}\right) + 2\cos^2 \frac{x}{2}, x \in \R \).

\begin{enumerate}
\item 求 \( f(x) \) 的值域;
\item 记 \(\triangle ABC\) 的内角 \(A, B, C\) 的对边长分别为 \(a, b, c\). 若 \(f(B) = 1\)，\(b = 1, c = \sqrt{3}\)，求 \(a\) 的值.
\end{enumerate}
\end{problem}

\begin{problem}
在甲，乙等 6 个单位参加的一次“唱读讲传”演出活动中，每个单位的节目集中安排在一起，若采用抽签的方式随机确定各单位的演出顺序（序号为 \(1, 2, \cdots, 6\)），求:

\begin{enumerate}
\item 甲，乙两单位的演出序号至少有一个为奇数的概率;
\item 甲，乙两单位之间的演出单位个数 \(\xi\) 的分布列与期望.
\end{enumerate}
\end{problem}

\begin{problem}
已知函数 \( f(x) = \frac{x - 1}{x + a} + \ln (x + 1) \)，其中实数 \( a \neq 1 \).

\begin{enumerate}
\item 若 \(a = 2\)，求曲线 \(y = f(x)\) 在点 \((0, f(0))\) 处的切线方程;
\item 若 \( f(x) \) 在 \( x = 1 \) 处取得极值，试讨论 \( f(x) \) 的单调性.
\end{enumerate}
\end{problem}

\begin{problem}
如图，四棱锥 \(P - ABCD\) 中，底面 \(ABCD\) 为矩形，\(PA \perp\) 底面 \(ABCD\)，\(PA = AB = \sqrt{6}\)，点 \(E\) 是棱 \(PB\) 的中点.

\begin{enumerate}
\item 求直线 AD 与平面 PBC 的距离;
\item 若 \(AD = \sqrt{3}\)，求二面角 \(A - EC - D\) 的平面角的余弦值.
\end{enumerate}

\bitmapfigure[max width=0.36\linewidth,max totalheight=4.8cm]{img/2010/chongqing_science/q19_fig1.png}

\end{problem}

\begin{problem}
已知以原点 \(O\) 为中心，\(F(\sqrt{5}, 0)\) 为右焦点的双曲线 \(C\) 的离心率 \(e = \frac{\sqrt{5}}{2}\).

\begin{enumerate}
\item 求双曲线 \(C\) 的标准方程及其渐近线方程;
\item 如图，已知过点 \(M(x_{1}, y_{1})\) 的直线 \(l_{1}: x_{1}x + 4y_{1}y = 4\) 与过点 \(N(x_{2}, y_{2})\)（其中 \(x_{2} \neq x_{1}\)）的直线 \(l_{2}: x_{2}x + 4y_{2}y = 4\) 的交点 \(E\) 在双曲线 \(C\) 上，直线 \(MN\) 与两条渐近线分别交于 \(G, H\) 两点，求 \(\triangle OGH\) 的面积.
\end{enumerate}

\bitmapfigure[max width=0.38\linewidth,max totalheight=4.8cm]{img/2010/chongqing_science/q20_fig1.png}

\end{problem}

\begin{problem}
在数列 \(\{a_{n}\}\) 中，\(a_{1} = 1\)，\(a_{n + 1} = c a_{n} + c^{n + 1}(2n + 1)(n \in \mathbf{N}^{*})\)，其中实数 \(c \neq 0\).

\begin{enumerate}
\item 求 \( \{a_{n}\} \) 的通项公式;
\item 若对一切 \(k \in \mathbf{N}^*\) 有 \(a_{2k} > a_{2k-1}\)，求 \(c\) 的取值范围.
\end{enumerate}
\end{problem}
